\paragrafos{Παράλληλες ευθείες και τέμνουσα}{Β.2.6, «Παράλληλες ευθείες που τέμνονται από μια άλλη ευθεία»}

\begin{anagke}
Ο δρόμος που περνά μπροστά από το σχολείο έχει δύο παράλληλα πεζοδρόμια. Ένα σοκάκι τον διασχίζει λοξά. Ο Νίκος παρατηρεί ότι στο ένα πεζοδρόμιο η γωνία μεταξύ σοκακιού και δρόμου είναι πολύ μυτερή, και ρωτάει: «Στο απέναντι πεζοδρόμιο θα είναι ίδια;». Ο Βασίλης αποφασίζει να τη μετρήσει: $50^\circ$ και στα δύο.

Η Αγγελική σχεδιάζει δύο παράλληλες ευθείες και μία τρίτη που τις τέμνει, και βρίσκει $8$ γωνίες. Η Πολυξένη λέει: «Αν ξέρω μία, ξέρω όλες!». Αλλά για να το πει με ακρίβεια, χρειαζόμαστε ονόματα για τις γωνίες: «εντός», «εκτός», «εναλλάξ», «επί τα αυτά».
\end{anagke}

\begin{lexi}{εντός, εκτός, εναλλάξ, επί τα αυτά}{θέσεις γωνιών}
Εντός σημαίνει «μέσα» στη λωρίδα των παραλλήλων, εκτός «έξω». Οι γωνίες \lexiword{εναλλάξ} είναι από αντίθετες πλευρές της τέμνουσας, ενώ οι γωνίες \lexiword{επί τα αυτά} είναι από την ίδια πλευρά της (το «επί τα αυτά μέρη» = προς τα ίδια μέρη).
\end{lexi}

\begin{metafrasi}
Στη λωρίδα ανάμεσα στις παράλληλες. & εντός\\
Έξω από τη λωρίδα. & εκτός\\
Από αντίθετες πλευρές της τέμνουσας. & εναλλάξ\\
Από την ίδια πλευρά της τέμνουσας. & επί τα αυτά\\
\end{metafrasi}

\tmima{Η θεωρία}

Έστω ευθείες $\varepsilon_1\parallel\varepsilon_2$ και μια \textbf{τέμνουσα} $\delta$ που τις τέμνει στα $\mathrm{Α}$ και $\mathrm{Β}$. Σχηματίζονται $8$ γωνίες· ονομάζουμε $\hat{\mathrm{Α}}_1,\hat{\mathrm{Α}}_2,\hat{\mathrm{Α}}_3,\hat{\mathrm{Α}}_4$ γύρω από το $\mathrm{Α}$ και $\hat{\mathrm{Β}}_1,\ldots,\hat{\mathrm{Β}}_4$ γύρω από το $\mathrm{Β}$.

\begin{center}
\begin{tikzpicture}[scale=1.0]
\draw[thick] (-2.6,1.2)--(2.6,1.2) node[right] {$\varepsilon_1$};
\draw[thick] (-2.6,-1.2)--(2.6,-1.2) node[right] {$\varepsilon_2$};
\draw[thick,gblue] (-1.2,2.4)--(1.2,-2.4) node[below right] {$\delta$};
\fill (-0.6,1.2) circle (2pt) node[above left] {$\mathrm{Α}$};
\fill (0.6,-1.2) circle (2pt) node[below right] {$\mathrm{Β}$};
\node at (-1.0,1.55) {\small $1$}; \node at (-0.2,1.55) {\small $2$};
\node at (-0.2,0.85) {\small $3$}; \node at (-1.0,0.85) {\small $4$};
\node at (0.2,-0.85) {\small $1$}; \node at (1.0,-0.85) {\small $2$};
\node at (1.0,-1.55) {\small $3$}; \node at (0.2,-1.55) {\small $4$};
\end{tikzpicture}
\end{center}

\begin{orismos}{ονομασίες γωνιών}
Οι γωνίες $\hat{\mathrm{Α}}_3,\hat{\mathrm{Α}}_4,\hat{\mathrm{Β}}_1,\hat{\mathrm{Β}}_2$ είναι \textbf{εντός}, οι υπόλοιπες \textbf{εκτός}. Δύο γωνίες είναι \textbf{εναλλάξ} αν βρίσκονται σε διαφορετικά ημιεπίπεδα της $\delta$ και \textbf{επί τα αυτά} αν βρίσκονται στο ίδιο. Προκύπτουν $6$ ονομασίες για τα $16$ ζεύγη γωνιών: εντός εναλλάξ, εκτός εναλλάξ, εντός και επί τα αυτά, εκτός και επί τα αυτά, εντός--εκτός εναλλάξ, εντός--εκτός επί τα αυτά.
\end{orismos}

\begin{kanonas}{παράλληλες και τέμνουσα}
Αν $\varepsilon_1\parallel\varepsilon_2$, τότε:
\begin{itemize}[leftmargin=1.2em,itemsep=0.3ex,topsep=2pt]
\item οι \textbf{εντός εναλλάξ} είναι ίσες,
\item οι \textbf{εκτός εναλλάξ} είναι ίσες,
\item οι \textbf{εντός--εκτός επί τα αυτά} είναι ίσες,
\item οι \textbf{εντός επί τα αυτά} είναι παραπληρωματικές ($\text{άθροισμα}=180^\circ$),
\item οι \textbf{εκτός επί τα αυτά} είναι παραπληρωματικές.
\end{itemize}
Δηλαδή όλες οι οξείες γωνίες του σχήματος είναι ίσες ($\omega$), όλες οι αμβλείες ίσες ($\varphi$) και $\omega+\varphi=180^\circ$.
\end{kanonas}

Ισχύει και το \textbf{αντίστροφο}: αν δύο ευθείες που τέμνονται από τρίτη σχηματίζουν εντός εναλλάξ γωνίες ίσες (ή εντός επί τα αυτά παραπληρωματικές, ή εντός--εκτός επί τα αυτά ίσες), τότε είναι παράλληλες. Έτσι ελέγχουμε αν δύο ευθείες είναι παράλληλες.

\begin{prosoxi}
\begin{itemize}[leftmargin=1.2em,itemsep=0.3ex,topsep=2pt]
\item Οι ιδιότητες ισχύουν \emph{μόνο όταν} οι ευθείες είναι παράλληλες.
\item Οι εντός επί τα αυτά είναι παραπληρωματικές (όχι ίσες).
\item Οι κατακορυφήν γωνίες είναι ίσες και οι διαδοχικές στο ίδιο σημείο παραπληρωματικές: $\hat{\mathrm{Α}}_1+\hat{\mathrm{Α}}_2=180^\circ$.
\end{itemize}
\end{prosoxi}

\tmima{Λυμένα παραδείγματα}

\paradeigma{Όλες οι γωνίες}
\tixero $\varepsilon_1\parallel\varepsilon_2$ και $\hat{\mathrm{Α}}_1=40^\circ$.\par
\tithelo Όλες οι άλλες γωνίες.\par
\possindeo Οι $\hat{\mathrm{Α}}_1$ και $\hat{\mathrm{Α}}_3$ είναι κατακορυφήν, άρα $\hat{\mathrm{Α}}_3=40^\circ$. Η $\hat{\mathrm{Α}}_2=180^\circ-40^\circ=140^\circ$ και η $\hat{\mathrm{Α}}_4=140^\circ$. Στο $\mathrm{Β}$ οι γωνίες είναι ίσες με τις αντίστοιχες του $\mathrm{Α}$ (εντός--εκτός επί τα αυτά).\par
\apantisi Σε κάθε σημείο: δύο γωνίες $40^\circ$ και δύο $140^\circ$.\par
\elegxos $40^\circ+140^\circ=180^\circ$ και το άθροισμα των γωνιών γύρω από σημείο είναι $2\cdot40^\circ+2\cdot140^\circ=360^\circ$.

\paradeigma{Εντός επί τα αυτά}
\tixero $\varepsilon_1\parallel\varepsilon_2$ και μία εντός επί τα αυτά γωνία είναι $65^\circ$.\par
\tithelo Η άλλη.\par
\possindeo Οι εντός επί τα αυτά είναι παραπληρωματικές: $x=180^\circ-65^\circ=115^\circ$.\par
\apantisi $115^\circ$.\par
\elegxos $65^\circ+115^\circ=180^\circ$.

\paradeigma{Παράλληλες ή όχι}
\tixero Δύο ευθείες που τέμνονται από τρίτη σχηματίζουν εντός εναλλάξ γωνίες $72^\circ$ και $75^\circ$.\par
\tithelo Είναι παράλληλες;\par
\possindeo Για να είναι παράλληλες θα έπρεπε οι εντός εναλλάξ να είναι ίσες.\par
\apantisi Όχι, γιατί $72^\circ\neq75^\circ$.\par
\elegxos Αν ήταν $72^\circ$ και $72^\circ$, θα ήταν παράλληλες.

\tmima{Ασκήσεις}

\zoni{Το κατάλαβα}{ονόματα και ισότητες}

\ex Αν $\varepsilon_1\parallel\varepsilon_2$ και μια οξεία γωνία είναι $55^\circ$, πόσες από τις $8$ γωνίες είναι $55^\circ$ και πόσες $125^\circ$;

\ex Αν δύο εντός εναλλάξ γωνίες είναι $x$ και $2x-30^\circ$ και οι ευθείες είναι παράλληλες, βρες το $x$.

\ex Δύο εντός επί τα αυτά γωνίες παραλλήλων είναι $x$ και $x+40^\circ$. Βρες τες.

\ex Δύο εκτός εναλλάξ γωνίες είναι $3x$ και $x+50^\circ$ και οι ευθείες παράλληλες. Βρες το $x$.

\ex Για κάθε ζεύγος γωνιών, πες πώς ονομάζεται: α) $\hat{\mathrm{Α}}_3,\hat{\mathrm{Β}}_1$ β) $\hat{\mathrm{Α}}_3,\hat{\mathrm{Β}}_2$ γ) $\hat{\mathrm{Α}}_1,\hat{\mathrm{Β}}_3$ δ) $\hat{\mathrm{Α}}_1,\hat{\mathrm{Β}}_1$.

\ex Σχεδίασε δύο παράλληλες που απέχουν $4$ cm και μια τέμνουσα που σχηματίζει $72^\circ$ με τη μία. Υπολόγισε τις υπόλοιπες γωνίες.

\ex Σωστό ή λάθος; α) Οι εντός εναλλάξ γωνίες είναι πάντα ίσες. \quad β) Αν δύο ευθείες έχουν εντός εναλλάξ γωνίες ίσες, είναι παράλληλες. \quad γ) Οι εντός επί τα αυτά παραλλήλων είναι ίσες.

\zoni{Το χρησιμοποιώ}{δρόμοι και σκάλες}

\ex \textbf{Το σοκάκι.} Στο πεζοδρόμιο η γωνία με το σοκάκι είναι $50^\circ$. Πόση είναι στο απέναντι πεζοδρόμιο (εντός εναλλάξ) και πόση η γωνία που είναι εντός επί τα αυτά με αυτήν;

\ex \textbf{Η σκάλα.} Μια σκάλα ακουμπά σε δύο παράλληλα ράφια. Η γωνία της με το πρώτο ράφι είναι $70^\circ$. Πόση είναι με το δεύτερο, από την ίδια πλευρά;

\ex \textbf{Οι ράγες.} Δύο παράλληλες σιδηροτροχιές τέμνονται από μια διάβαση. Αν μια γωνία είναι $110^\circ$, βρες όλες τις άλλες.

\ex \textbf{Ο φράκτης.} Δύο πασσάλοι κεκλιμένοι προς το έδαφος σχηματίζουν γωνίες $64^\circ$ και $64^\circ$ (εντός εναλλάξ) με το σχοινί. Είναι παράλληλοι; Γιατί;

\ex \textbf{Το ράφι.} Δύο ράφια τέμνονται από κατακόρυφο στύλο και σχηματίζουν εντός επί τα αυτά γωνίες $90^\circ$ και $90^\circ$. Τι ισχύει για τα ράφια;

\zoni{Το σκέφτομαι}{αποδείξεις}

\ex Εξήγησε γιατί αν οι εντός εναλλάξ είναι ίσες τότε και οι εντός επί τα αυτά είναι παραπληρωματικές.

\ex Στο τρίγωνο $\mathrm{ΑΒΓ}$ φέρε από το $\mathrm{Α}$ παράλληλη προς τη $\mathrm{ΒΓ}$. Τι ισχύει για τις γωνίες που σχηματίζονται;

\ex Αν $\varepsilon_1\parallel\varepsilon_2$ και $\varepsilon_3\parallel\varepsilon_4$, και μια γωνία τους είναι $70^\circ$, βρες όλες τις γωνίες του «παραλληλογράμμου» που σχηματίζεται.

\ex Αν η $\mathrm{Βδ}$ είναι διχοτόμος της γωνίας $\hat{\mathrm{Β}}=56^\circ$ και $\varepsilon_1\parallel\varepsilon_2$, βρες τη γωνία που σχηματίζει η διχοτόμος με την $\varepsilon_1$.

\zoni{Για το σπίτι}{ρίξε μια ματιά}

\ex Βρες στο σπίτι δύο παράλληλες ευθείες τεμνόμενες από τρίτη και μέτρησε τις γωνίες.

\ex Σχεδίασε τετράπλευρο $\mathrm{ΑΒΓΔ}$ με $\mathrm{ΑΒ}\parallel\mathrm{ΔΓ}$, $\mathrm{ΑΔ}\parallel\mathrm{ΒΓ}$ και $\hat{\mathrm{Α}}=63^\circ$. Βρες τις άλλες γωνίες.
